\chapter{Introduzione}
\label{introduzione}

L'evoluzione delle reti mobili verso la quinta generazione ha introdotto un'architettura più articolata, composta da diverse funzioni di rete che cooperano per erogare i servizi del sistema. In questo contesto, la verifica della sicurezza assume un ruolo importante ed è supportata dalle \emph{Security Assurance Specifications} (SCAS), che definiscono requisiti e procedure di test per valutare specifiche proprietà di sicurezza \cite{3gpp_ts33117_2019,gsma_fs13_2025}. Tuttavia, determinare l'esito di una prova a partire dalle catture di traffico non è sempre immediato: i criteri di verifica sono descritti in linguaggio naturale, mentre le evidenze necessarie possono essere distribuite tra messaggi e protocolli differenti.

Il presente lavoro affronta questo problema proponendo un \emph{Domain-Specific Language} (DSL) che permette di esprimere i criteri di verifica in forma dichiarativa, separandoli dalla logica del software che li valuta. Il sistema sviluppato seleziona gli eventi rilevanti dalle catture di traffico, applica le regole specificate e produce un report contenente il verdetto e le evidenze che lo giustificano. La soluzione è stata valutata su diversi casi SCAS e attraverso esperimenti di \emph{mutation testing}, nei quali vengono introdotte alterazioni controllate nelle catture per verificare la capacità del sistema di rilevare le modifiche apportate.
